\SGANextSectionTarget{sga3:xix:section:5}\SGASection{5}{An instructive example}

\par\medskip\noindent
\textbf{\SGAInlineRefTarget{sga3:xix:paragraph:5-1}5.1.}---
Let \(k\) be an algebraically closed field of characteristic \(0\).
Put \(A=k[t]\), the polynomial ring in one variable over \(k\), and
let \(S=\operatorname{Spec}(A)\).  Consider the following Lie algebra
\(\mathfrak g\) over \(\mathcal O_S\).  As a module it is free of
rank \(3\), with basis \(\{X,Y,H\}\), and its multiplication table is
\[
  [X,Y]=2tH,
  \qquad
  [H,X]=X,
  \qquad
  [H,Y]=-Y.
\]
For \(s\in S\), \(s\ne s_0\), where \(s_0\) is the point defined by
\(t=0\), the fiber
\[
  \mathfrak g(s)=\mathfrak g\otimes_A\kappa(s)
\]
is isomorphic to the Lie algebra of
\(\mathrm{PGL}_{2,\kappa(s)}\).  For \(s=s_0\), it is a solvable Lie
algebra.
\SGARefTarget{sga3:xix:paragraph:5:1}

\par\medskip\noindent
\textbf{\SGAInlineRefTarget{sga3:xix:paragraph:5-2}5.2.}---
Let \(G_1\) be the group scheme of automorphisms of \(\mathfrak g\):
for every \(S'\to S\), \(G_1(S')\) is the group of automorphisms of
the \(\mathcal O_{S'}\)-Lie algebra
\(\mathfrak g\otimes_{\mathcal O_S}\mathcal O_{S'}\).  It is a closed
subscheme of the group \(\mathrm{GL}(\mathfrak g)\) of automorphisms
of the \(\mathcal O_S\)-module \(\mathfrak g\).

Let \(S'\to S\), and let
\[
  u\in M_3\bigl(\Gamma(S',\mathcal O_{S'})\bigr)
\]
be viewed as an endomorphism of the \(\mathcal O_{S'}\)-module
\(\mathfrak g\otimes_{\mathcal O_S}\mathcal O_{S'}\):
\[
  \begin{aligned}
  u(X)&=aX+bY+eH,\\
  u(Y)&=b'X+a'Y+e'H,\\
  u(H)&=cX+c'Y+dH.
  \end{aligned}
\]
It is immediate that \(u\) is a section of \(G_1\) if and only if
\(\det(u)\) is invertible and the following relations hold:
\[
\begin{alignedat}{4}
  \textup{(1)}\quad&a(d-1)=ec,
  &\qquad&
  \textup{(1')}\quad&a'(d-1)=e'c',\\
  \textup{(2)}\quad&b(d+1)=ec',
  &&
  \textup{(2')}\quad&b'(d+1)=e'c,\\
  \textup{(3)}\quad&e=2t(bc-ac'),
  &&
  \textup{(3')}\quad&e'=2t(b'c'-a'c),\\
  \textup{(4)}\quad&2tc=eb'-ae',
  &&
  \textup{(4')}\quad&2tc'=be'-ea',\\
  \textup{(5)}\quad&2t(aa'-bb')=2td.&&&
\end{alignedat}
\]
\SGARefTarget{sga3:xix:paragraph:5:2}

\par\medskip\noindent
\SGAInlineRefTarget{sga3:xix:lemma:5-3}\textbf{\textit{\SGARefLink{sga3:xix:lemma:5:3}{Lemma 5.3}.}}---
\textit{Relations \textup{(1)}, \textup{(1')}, \textup{(2)}, and
\textup{(2')} imply}
\[
  \det(u)=aa'(2-d)+bb'(d+2),
  \qquad
  aa'-bb'=d\det(u).
\]
\SGARefTarget{sga3:xix:lemma:5:3}

Indeed, the first assertion follows at once by substituting
relations \textup{(1)}, \textup{(1')}, \textup{(2)}, and
\textup{(2')} in the expansion of \(\det(u)\):
\[
\begin{aligned}
  \det(u)
  &=aa'd+be'c+b'c'e-a'ec-ae'c'-bb'd\\
  &=aa'd+bb'(d+1)+bb'(d+1)
    -aa'(d-1)-bb'd-aa'(d-1)\\
  &=aa'(d-d+1-d+1)+bb'(d+1+d+1-d)\\
  &=aa'(2-d)+bb'(d+2).
\end{aligned}
\]
Multiplying this relation by \(d\), we obtain
\[
  d\det(u)=aa'(2d-d^2)+bb'(d^2+2d).
\]

